% ============================================================
% SigmaFlow v8 — Editable LaTeX Template (main.tex)
% ============================================================
%
% This template is the base structure used by:
%     sigmaflow/report/latex_report.py
%
% You can MANUALLY EDIT this file to customize the report.
% Placeholders ({{...}}) are filled automatically by the engine.
%
% To compile manually:
%   pdflatex main.tex
%   pdflatex main.tex    (second pass for TOC)
% ============================================================

\documentclass[12pt, a4paper]{article}

\usepackage[a4paper, top=2.5cm, bottom=2.5cm, left=2.5cm, right=2.5cm]{geometry}
\usepackage{graphicx, float, booktabs, longtable, hyperref, xcolor}
\usepackage{enumitem, fancyhdr, titlesec, array, amsmath, microtype}
\usepackage[T1]{fontenc}
\usepackage[utf8]{inputenc}
\usepackage{lmodern}

\definecolor{sigmaBlue}{RGB}{21, 101, 192}
\definecolor{sigmaGreen}{RGB}{46, 125, 50}
\definecolor{sigmaRed}{RGB}{198, 40, 40}
\definecolor{sigmaOrange}{RGB}{230, 81, 0}
\definecolor{sigmaGray}{RGB}{117, 117, 117}
\definecolor{abstractBg}{RGB}{240, 244, 255}

\hypersetup{colorlinks=true, linkcolor=sigmaBlue, urlcolor=sigmaBlue}

\titleformat{\section}{\color{sigmaBlue}\Large\bfseries}{\thesection}{1em}{}
\titleformat{\subsection}{\color{sigmaBlue}\large\bfseries}{\thesubsection}{1em}{}

\pagestyle{fancy}
\fancyhf{}
\lhead{\small\color{sigmaGray}SigmaFlow v8 — Automated DMAIC Report}
\rhead{\small\color{sigmaGray}\today}
\cfoot{\small\thepage}

\title{
  \vspace{1.5cm}
  {\color{sigmaBlue}\Huge\bfseries SigmaFlow v8\\[0.3em]
  Automated Statistical Analysis Report}\\[0.7em]
  {\large\color{sigmaGray} Lean Six Sigma / DMAIC Framework}
}
\author{}
\date{\today}

\begin{document}
\maketitle
\thispagestyle{empty}
\newpage

% ─── Abstract ───────────────────────────────────────────────
\section*{Abstract}
\addcontentsline{toc}{section}{Abstract}

\begin{flushleft}
\colorbox{abstractBg}{\parbox{0.96\textwidth}{
\smallskip
{{abstract_text}}
\smallskip
}}
\end{flushleft}

\newpage
\tableofcontents
\newpage

% ─── Introduction ───────────────────────────────────────────
\section{Introduction}

This report was automatically generated by \textbf{SigmaFlow v8}, a modular
Python framework for industrial Six Sigma analysis. The analysis pipeline
detects dataset type, runs statistical analysis, generates visualizations,
applies Western Electric rules, performs root cause correlation analysis,
and compiles all findings into this document.

% ─── Methodology ────────────────────────────────────────────
\section{Methodology}

\subsection{Statistical Process Control (SPC)}
Control charts distinguish common-cause (random) from special-cause
(assignable) variation. The XmR chart applies $\pm 3\sigma$ limits to
individual observations.

\subsection{Western Electric Rules}
\begin{itemize}
  \item \textbf{Rule 1} — Any point outside $\pm 3\sigma$ ($p \approx 0.27\%$)
  \item \textbf{Rule 2} — Nine consecutive points on one side of the center line
  \item \textbf{Rule 3} — Six consecutive trending points (drift indicator)
  \item \textbf{Rule 4} — Fourteen alternating points (tampering indicator)
\end{itemize}

\subsection{Process Capability}
\begin{center}
\begin{tabular}{ll}
\toprule \textbf{Index} & \textbf{Interpretation threshold} \\
\midrule
$C_{pk} < 1.00$ & Not capable — immediate action \\
$1.00 \le C_{pk} < 1.33$ & Marginal — monitor closely \\
$1.33 \le C_{pk} < 1.67$ & Acceptable \\
$C_{pk} \ge 1.67$ & Excellent (Six Sigma level) \\
\bottomrule
\end{tabular}
\end{center}

\subsection{Root Cause Analysis}
Pearson and Spearman correlations identify variables most associated
with the quality target. $|r| \ge 0.70$ = strong candidate for investigation.

% ─── Dataset Description ────────────────────────────────────
\section{Dataset Description}

{{dataset_description}}

% ─── Statistical Analysis ───────────────────────────────────
\section{Statistical Analysis}

{{analysis_summary}}

% ─── Process Stability ──────────────────────────────────────
\section{Process Stability Analysis}

\begin{figure}[H]
  \centering
  \includegraphics[width=0.92\textwidth]{figures/spc_xmr_chart.png}
  \caption{XmR Control Chart — Individuals and Moving Range}
\end{figure}

% ─── Capability Analysis ────────────────────────────────────
\section{Capability Analysis}

{{capability_section}}

\begin{figure}[H]
  \centering
  \includegraphics[width=0.92\textwidth]{figures/capability_plot.png}
  \caption{Process Capability Plot}
\end{figure}

% ─── Detected Anomalies ─────────────────────────────────────
\section{Detected Anomalies}

{{anomalies_section}}

% ─── Example manually-written anomaly block ─────────────────
\subsection*{Western Electric Rule 1 — Interpretation Guide}

A point exceeding the $3\sigma$ control limit was detected.

\textbf{Interpretation:} This indicates a special cause of variation.
A point beyond $\pm 3\sigma$ has only a 0.27\% probability under a
stable normal process.

\textbf{Possible Causes:}
\begin{itemize}
  \item Equipment malfunction or sudden parameter shift
  \item Raw material variation (new batch / supplier)
  \item Operator change or training gap
  \item Measurement system drift or calibration error
\end{itemize}

\textbf{Recommended Investigation:}
Inspect process logs at the flagged timestamp. Check calibration records
and maintenance logs. Apply the 5-Why method to trace to root cause.

% ─── Root Cause Indicators ──────────────────────────────────
\section{Root Cause Indicators}

{{root_cause_section}}

\begin{figure}[H]
  \centering
  \includegraphics[width=0.92\textwidth]{figures/correlation_heatmap.png}
  \caption{Correlation Matrix — Root Cause Analysis}
\end{figure}

\begin{figure}[H]
  \centering
  \includegraphics[width=0.92\textwidth]{figures/variable_importance.png}
  \caption{Variable Importance Ranking}
\end{figure}

% ─── Recommendations ────────────────────────────────────────
\section{Recommendations}

{{recommendations}}

% ─── Conclusion ─────────────────────────────────────────────
\section{Conclusion}

{{conclusion}}

\end{document}
